\scanpage{113}
\diagram{\begin{tikzcd}[ampersand replacement=\&,column sep=.55em,row sep=.65em]
 \&\deleted{\unsure{G'_1}\ P_1\ G_1}\arrow[dl,dash]\arrow[dr,dash]\&\&\&P_1\arrow[dl,dash,"G_2"']\arrow[dr,dash,"G_3"]\&\\
 \deleted{G_2\ P_2}\arrow[rr,dash,"\deleted{\unsure{G_2\,G_3}}"']\&\&\deleted{P_3\,G'_3}\&
 P_3\arrow[rr,dash,"G_1"']\&\&P_2
\end{tikzcd}}
\[
\begin{array}{lll@{\qquad}lll}
 P_1&\text{bitorseur sous}&G_2,G_3&P_1^{\mathrm o}&\text{bitorseur sous}&G_3,G_2\\
 P_2&\text{-------------}&G_3,G_1&P_2^{\mathrm o}&\text{---------}&G_1,G_3\\
 P_3&\text{-------------}&G_1,G_2&P_3^{\mathrm o}&\text{---------}&G_2,G_1
\end{array}
\]
\margin{\(P_1^{\mathrm o}\nearrow\) opposé de \(P_1\)}
\[
 P_1\wedge P_2\wedge P_3\quad\text{bitorseur sous }G_2,G_2\quad\text{etc.}\ /\!/\quad
 P_3^{\mathrm o}\wedge P_2^{\mathrm o}\wedge P_1^{\mathrm o}\quad\text{\unsure{bit. ss|bitors. ss|bitors. sous}}\,G_{\unsure{1|2}},G_2
\]
\[
\begin{aligned}
 P_1\wedge P_2\wedge P_3&\xrightarrow[\sim]{\varphi_2}\deleted{\unsure{G_2}}{}_{G_2}\mathbb{1}_{G_2}
 &P_3^{\mathrm o}\wedge P_2^{\mathrm o}\wedge P_1^{\mathrm o}&\overset{\varphi_2^{\mathrm o}}{\simeq}\deleted{\unsure{\operatorname{bit}}}{}_{G_2}\mathbb{1}_{G_2}\\
 P_2\wedge P_3\wedge P_1&\overset{\varphi_3}{\simeq}{}_{G_3}\mathbb{1}_{G_3}
 \quad\deleted{\unsure{\mathbb{1}|1|\Updownarrow}}&\deleted{\unsure{P}}P_1^{\mathrm o}\wedge P_2^{\mathrm o}\wedge P_3^{\mathrm o}&\overset{\varphi_3^{\mathrm o}}{\simeq}{}_{G_3}\mathbb{1}_{G_3}\\
 P_3\wedge P_1\wedge P_2&\overset{\varphi_1}{\simeq}{}_{G_1}\mathbb{1}_{G_1}
 &P_2^{\mathrm o}\wedge P_1^{\mathrm o}\wedge P_3^{\mathrm o}&\overset{\varphi_1^{\mathrm o}}{\simeq}{}_{G_1}\mathbb{1}_{G_1}
\end{aligned}
\]
\[
 \Gamma\subset P_1\times P_2\times P_3\quad\inserted{\Gamma\ne\varnothing}\quad\text{tel que}
\]
\[
 \left\{\begin{aligned}
 x=(x_1,x_2,x_3)&\in P_1\times P_2\times P_3\\
 g_1\in G_1,\ g_2&\in G_2,\ g_3\in G_3
 \end{aligned}\right.
 \Rightarrow\left(\begin{aligned}
 &(x_1,x_2,x_3)\in \Gamma\\
 &\Longleftrightarrow g_2x_1g_3^{-1},g_3x_2g_1^{-1},g_1x_3g_2^{-1}\\
 &\hspace{12em}\in \Gamma
 \end{aligned}\right)
 \tag*{(A)}
\]
NB \(\ G_1\times G_2\times G_3\) \unsure{opère|opèrent} sur \(P_1\times P_2\times P_3\) par
\[
 \sigma(g_1,g_2,g_3)\mathbin{:}(x_1,x_2,x_3)
 =(g_2x_1g_3^{-1},g_3x_2g_1^{-1},g_1x_3g_2^{-1})
\]
et l'axiome envisagé, c'est que \(\Gamma\) est une orbite sous \(G_1\times G_2\times G_3\)
\[
 \Gamma\xrightarrow{p_i}P_i\quad\text{induite par \unsure{pr.|pr}$_i$}
\]
\deleted{\unsure{que}} Supposons fixé un pt \(\gamma\in \Gamma\), d'où des pts \(a_i=p_i(\gamma)\), d'où des iso.
\[
\begin{aligned}
 G_2&\xrightarrow{\widetilde a_1}G_3&
 &g_2\text{ et }g_3\text{ se correspondent}\\
 &&&\text{ssi}\quad g_2a_1=a_1g_3\\
 G_3&\overset{\widetilde a_2}{\relbar\joinrel\relbar} G_1&
 &g_3\text{ et }g_1\quad\text{---}\qquad g_3a_2=a_2g_1\\
 G_1&\overset{\widetilde a_3}{\longleftrightarrow}G_2&
 &g_1\text{ et }g_2\quad\text{---}\qquad g_1a_3=a_3g_2^{\deleted{\unsure{-1}}}
\end{aligned}
\]
C'est un syst. transitif d'iso, \unsure{donc|car|i.e.}
% archive sheet number: 114
